% ECONOMICS_PROBLEM: {"id": "D72-545", "title": "Persistence of clientelism", "jel_code": "D72", "source_id": "OP-0126"}
% !TeX program = xelatex
\documentclass[11pt,a4paper]{article}
\usepackage[margin=23mm]{geometry}
\usepackage{fontspec}
\setmainfont{Latin Modern Roman}
\usepackage{amsmath,amssymb,mathtools}
\usepackage{xcolor,enumitem,booktabs,longtable,array,xurl,hyperref}
\definecolor{muted}{HTML}{53636C}
\hypersetup{colorlinks=true,urlcolor=blue,linkcolor=blue}
\setlength{\parindent}{0pt}
\setlength{\parskip}{5pt}
\newcommand{\R}{\mathbb R}
\newcommand{\E}{\mathbb E}
\newcommand{\N}{\mathbb N}
\newcommand{\ind}{\mathbf 1}
\newcommand{\Var}{\operatorname{Var}}
\newcommand{\Cov}{\operatorname{Cov}}
\newcommand{\supp}{\operatorname{supp}}
\DeclareMathOperator*{\argmin}{arg\,min}
\DeclareMathOperator*{\argmax}{arg\,max}
\newcommand{\entrymeta}[3]{{\sffamily\small\color{muted}\textbf{\texttt{#1}}\quad|\quad #2\par #3\par}}
\newcommand{\Problemref}[1]{\texttt{#1}}
\newcommand{\JELref}[2]{\texttt{#2} (JEL \texttt{#1})}
\begin{document}
\section*{Persistence of clientelism}
\textbf{Problem ID:} \texttt{D72-545}\quad \textbf{JEL:} \texttt{D72}\par
% BEGIN ECONOMICS STATEMENT
\label{problem:OP-0126}
\label{atlas:prob:Q6}
{\sffamily\small\color{muted}\textbf{Q6} | Political economy | Editorial research problem family\par}
\subsection*{Definitions and assumptions}
Specify an infinite repeated game with a finite voter set $N$, parties and brokers, bounded per-period utilities, and discount factors $\delta_i\in(0,1)$. At each date, parties offer feasible benefits $b_{it}$, voter $i$ privately chooses support $v_{it}\in\{0,1\}$, and brokers receive signal $s_{it}$ with known conditional law $q(s\mid v,b,H_t)$. The public history $H_t$ contains only declared observable information. Voter utility includes material benefits, policy preferences and, if allowed, an explicitly specified reciprocity term. A policy $a$ changes ballot secrecy, broker information or access to programmatic services. Let $\mathcal E(m,a)$ be the set of sequential equilibria of the declared game. Require uniformly bounded transfers and choose an evaluation discount factor $\rho\in(0,1)$. For equilibrium $e$, define $B_m(e,a)=(1-\rho)\mathbb E_{m,e,a}[\sum_{t=0}^{\infty}\rho^t |N|^{-1}\sum_i b_{it}]$ and $V_m(e,a)=(1-\rho)\mathbb E_{m,e,a}[\sum_{t=0}^{\infty}\rho^t |N|^{-1}\sum_i v_{it}]$; the probability law depends on the model, equilibrium and policy.
\subsection*{Formal research question}
Characterize when $\mathcal E(m,a)$ contains equilibria with positive targeted transfers and a higher support probability than under the no-transfer comparison, and identify restrictions on signal informativeness, reciprocity and public-service provision that alter those outcomes. For policy evaluation, fix an equilibrium selector $\sigma(m,a)$ or report the set
\[
\{(B_m(e,a),V_m(e,a),W_m(e,a)):e\in\mathcal E(m,a)\},
\]
where $W_m$ uses declared welfare weights. Determine whether experimental and network data distinguish reciprocity from threats contingent on inferred voting.
\subsection*{Interpretation and scope}
Perfect ballot secrecy means that, conditional on the same prior history and offer, the broker's signal distribution does not depend on the individual vote. A sanction based only on such a signal cannot selectively punish a deviation in that vote. This does not exclude monitoring turnout, group outcomes, repeated relationships or reciprocity; those must be modeled separately. Persistence is therefore a conditional equilibrium and empirical question, not a contradiction between clientelism and secret ballots.
\subsection*{Source audit and status}
All three original references are verified. Stokes models imperfect inference, Wantchekon studies campaign appeals, and Finan--Schechter examines reciprocal targeting. They already provide mechanisms for persistence. The retained question concerns distinguishing and extending those mechanisms under specified reforms, rather than asking whether any explanation exists. The game and comparative-statics target are editorial specifications.
\subsection*{References}
{\footnotesize\raggedright
\begin{enumerate}[label={\textbf{[S\arabic*]}},leftmargin=3em]
\item Susan C. Stokes (2005). \emph{Perverse Accountability: A Formal Model of Machine Politics with Evidence from Argentina}. American Political Science Review 99(3), 315--325.\par
\textit{Locator and source role:} Abstract and repeated-game mechanism; imperfect inference about secret ballots rather than perfect observation of votes.\par
\url{https://doi.org/10.1017/S0003055405051683}
\item Leonard Wantchekon (2003). \emph{Clientelism and Voting Behavior: Evidence from a Field Experiment in Benin}. World Politics 55(3), 399--422.\par
\textit{Locator and source role:} Abstract and campaign experiment; constituency-service and patronage appeals, not simply verified purchases of individual secret votes.\par
\url{https://www.cambridge.org/core/journals/world-politics/article/abs/clientelism-and-voting-behavior-evidence-from-a-field-experiment-in-benin/3E386064D15E5E162AEDCEBECB32E8CB}
\item Frederico Finan and Laura Schechter (2012). \emph{Vote-Buying and Reciprocity}. Econometrica 80(2), 863--881.\par
\textit{Locator and source role:} Abstract and survey/experimental reciprocity measure; an alternative mechanism supporting transfers under secrecy.\par
\url{https://onlinelibrary.wiley.com/doi/10.3982/ECTA9035}
\end{enumerate}
}

\subsection*{Common conventions: Source mathematical conventions}

Unless an entry states otherwise, agent and firm sets are finite. State and action spaces are measurable, strategies and outcome maps are measurable, probability laws are specified, and expectations exist for the targets under discussion. Infinite-horizon discount factors lie in $(0,1)$ and discounted objectives are finite. Norms, tolerances and complexity input sizes must be specified for computational comparisons. Constants or primitives that appear locally are inputs, not empirical facts asserted by this catalogue.

\textbf{Identification.} A statistical model $m\in\mathcal M$ implies an observable law $P_m$ and target $\theta(m)$. At law $P$, its identified set is
\[
\Theta(P;\mathcal M)=\{\theta(m):m\in\mathcal M,\ P_m=P\}.
\]
Point identification holds when this set is a singleton, or equivalently when $P_{m_1}=P_{m_2}$ implies $\theta(m_1)=\theta(m_2)$ for all compatible models. Sharp bounds mean bounds attained, or approached when the boundary is not attained, by compatible models. Writing this set is a definition, not a solution: the substantive task is to establish informative restrictions and derive or estimate the set. An unrestricted-confounding model may give only trivial bounds.

\textbf{Inference.} Identification, estimability and valid uncertainty quantification are separate questions. Fix $\alpha\in(0,1)$. A confidence procedure $C_n$ over a declared law class $\mathcal P$ has asymptotic uniform coverage at level $1-\alpha$ when
\[
\liminf_{n\to\infty}\inf_{P\in\mathcal P}\Pr_P\{\theta(P)\in C_n\}\geq1-\alpha.
\]
For set-identified targets the coverage event must be defined separately, for example coverage of the full identified set. If an entry uses identification-indexed models instead of uniquely indexed laws, the same coverage convention is applied over those models. Pointwise limits do not establish uniform validity, and asymptotic validity does not guarantee finite-sample precision. Sample sizes, dependence structures and weak-identification sequences are part of each application.

\textbf{Counterfactuals.} $Y_i(a)$ is a potential outcome under a specified intervention, including its timing, implementation and versions. It is not $Y_i$ conditional on voluntarily choosing $a$. A full intervention vector permits interference; shorthand scalar treatment notation does not silently impose no interference. Assignment, selection, attrition, anticipation and measurement must be specified. A claim about an instrument's local effect is not automatically a population policy effect.

\textbf{Economic equilibrium.} A counterfactual model includes best responses, constraints, feasibility, market clearing, state transitions and expectations consistent with the stated information. When several equilibria exist, either a declared selector is an input or the target is set-valued. Comparisons across policy regimes must use the stated selection convention. Reduced-form relationships are not presumed invariant after an equilibrium-changing intervention.

\textbf{Optimization and welfare.} Optimization is over the stated feasible set. If attainment is not established, $\sup$ or $\inf$ and $\varepsilon$-optimal choices replace an asserted maximizer or minimizer. For example, with value $V=\sup_{a\in\mathcal A}W(a)<\infty$, an $\varepsilon$-optimal set is $\{a\in\mathcal A:W(a)\geq V-\varepsilon\}$ for $\varepsilon>0$. Benefits, costs, risk and health or environmental outcomes can be aggregated only using declared units and conversion weights. Interpersonal utility weights, the treatment of fiscal transfers and foreign profits, discounting, and the behavioral welfare criterion are explicit inputs. A comparative-static derivative additionally requires existence and appropriate differentiability of the selected counterfactual.

\textbf{Sources.} Local labels [S1], [S2], and so forth restart in every question. Locators identify the relevant abstract, model, section or result at the level actually checked; a bibliographic or abstract-level verification is not represented as inspection of an entire proof. Working papers are distinguished from later publications and changing versions. Source summaries are paraphrased. All URL checks and editorial formulations refer to the audit date stated above.
% END ECONOMICS STATEMENT
\end{document}
